% Public source snapshot of the instructor's Module 11 alignment lecture. Private
% links, executable implementations, diagram source, external images, paper
% screenshots, and product-current claims are omitted under sources/README.md.
% The preserved prose and equations document provenance; publication claims are
% independently checked against primary sources.
\documentclass[xcolor=svgnames, 10pt, handout]{beamer}
\usepackage{amsmath}
\usepackage{amssymb}
\usepackage{booktabs}
\usepackage{xcolor}
% Define colors for better visualization
\definecolor{darkblue}{rgb}{0.0, 0.0, 0.55}
\definecolor{darkred}{rgb}{0.55, 0.0, 0.0}

\title{Aligning LLMs with Human Preferences}
% \subtitle{RLHF, DPO, and the Science of Helpful and Harmless AI}
\author{Heman Shakeri}
\date{}

\begin{document}

\begin{frame}
    \titlepage
\end{frame}

\begin{frame}{Public Snapshot Boundary}
    \begin{alertblock}{Provenance, not a publishable asset bundle}
        This snapshot preserves the instructor-authored lecture spine and equations.
        Private links, executable or library code, diagram source, external images and
        paper screenshots, and claims tied to current products have been omitted; their
        original locations remain marked. Remaining claims are historical lecture notes
        and are independently checked against primary sources before book publication.
    \end{alertblock}
\end{frame}

\begin{frame}{Agenda}
    \tableofcontents
\end{frame}

% ----------------------------------------------------------------
\section{Introduction: The Alignment Problem}
% ----------------------------------------------------------------

\begin{frame}{Recap: Building Capable Models}
    We have established how to adapt and enhance Large Language Models:
    \begin{itemize}
        \item \textbf{Pre-training:} Building foundational knowledge.
        \pause
        \item \textbf{Supervised Fine-Tuning (SFT):} Adapting models for specific tasks or instruction following.
        \pause
        \item \textbf{PEFT (LoRA/QLoRA):} Efficiently performing SFT.
        \pause
        \item \textbf{RAG:} Connecting LLMs to external, dynamic data.
    \end{itemize}
    \vspace{0.5cm}
    \pause
    These techniques focus on \textit{capabilities}---what the model \textit{can} do.
    \pause

    Now we focus on \textit{behavior}---what the model \textit{should} do.
    \pause

    \begin{alertblock}{From Model to Assistant}
        Transforming a capable language model into a helpful, safe assistant requires a fundamentally different type of training.
    \end{alertblock}
\end{frame}

\begin{frame}{The Limitations of SFT}
    SFT maximizes the likelihood of the training data (MLE). It teaches the model the \textit{format} of interaction, but struggles to define the \textit{quality} of interaction.
    \pause

    \begin{block}{The Alignment Gap}
        A model can follow instructions correctly according to SFT loss, yet still be undesirable.
    \end{block}
    \pause

    We aim for models that are \textbf{HHH} (Anthropic's framework):
    \begin{itemize}
        \item \textbf{Helpful:} Provides useful, well-formatted, and detailed answers.
        \pause
        \item \textbf{Harmless:} Avoids generating toxic, biased, or dangerous content.
        \pause
        \item \textbf{Honest:} Provides accurate information and\\ expresses uncertainty (avoids hallucination).
    \end{itemize}
    \pause
    This challenge of instilling human values and\\ preferences is the \textbf{Alignment Problem}.

    \only<6->{
        % [Diagram or external-image composition omitted from the public snapshot; the book uses independently drawn figures.]
    }
\end{frame}

\begin{frame}{Why Preferences Instead of Demonstrations?}
    Why not just use SFT with ``perfect'' examples?
    \pause
    \begin{enumerate}
        \item \textbf{Difficulty of Specification:} It is hard to write explicit rules for subjective concepts like ``helpfulness'' or ``creativity.''
        \pause
        \item \textbf{Generation vs. Evaluation Gap:} Humans are much better at \textit{comparing} two generated outputs (A vs B) than \textit{writing} the ideal output from scratch.
        \pause
        \item \textbf{Scalability:} Collecting comparison data is often faster and provides a denser learning signal than expert demonstrations.
    \end{enumerate}
    \pause
    \begin{block}{The Core Idea}
        Leverage human feedback in the form of \textbf{preferences} to guide the optimization process beyond SFT.
    \end{block}
\end{frame}

% ----------------------------------------------------------------
\section{Background: Deep Reinforcement Learning}
% ----------------------------------------------------------------

\begin{frame}{Reinforcement Learning Primer}
    \textbf{Reinforcement Learning (RL)} is a paradigm where an agent learns to make decisions by interacting with an environment.
    \pause

    \begin{center}
    % [Diagram or external-image composition omitted from the public snapshot; the book uses independently drawn figures.]
    \end{center}
    \pause

    \textbf{Key Components:}
    \begin{itemize}
        \item \textbf{Agent:} The learner/decision maker (e.g., a neural network).
        \item \textbf{State ($s_t$):} Current situation/observation.
        \item \textbf{Action ($a_t$):} Decision made by the agent.
        \item \textbf{Reward ($r_t$):} Scalar feedback signal (how good was the action?).
        \item \textbf{Policy ($\pi$):} Agent's strategy, mapping states to actions.
    \end{itemize}
\end{frame}

\begin{frame}{RL for Language Models: The Connection}
    How do we map RL concepts to LLM training?
    \pause

    \begin{center}
    % [Diagram or external-image composition omitted from the public snapshot; the book uses independently drawn figures.]
    \end{center}
    \pause

    \textbf{Goal:} Train policy $\pi_\theta$ to maximize expected cumulative reward (i.e., generate high-quality, complete responses).
\end{frame}

\begin{frame}{Proximal Policy Optimization (PPO)}
    \textbf{PPO} (Schulman et al., 2017) is the RL algorithm developed in this lecture's classic alignment pipeline.
    \pause

    \textbf{Intuition:} Limit how much the policy can change \textit{per optimization step}.
    \pause
    \textbf{Why PPO?}
    \begin{itemize}
        \item Policy gradient methods can be unstable---large updates can catastrophically hurt performance.
        \pause
        \item PPO constrains updates to stay within a ``trust region'' of the \textit{previous iteration's policy} ($\pi_{\theta}^{\text{old}} \rightarrow \pi_{\theta}^{\text{new}}$).
        \pause
        \item More stable and sample-efficient than vanilla policy gradients.
    \end{itemize}
    \pause

    \textbf{Key Mechanism: Clipped Objective}
    \begin{equation*}
        L^{CLIP}(\theta) = \mathbb{E}_t \left[ \min\left( r_t(\theta) \hat{A}_t, \text{clip}(r_t(\theta), 1-\epsilon, 1+\epsilon) \hat{A}_t \right) \right]
    \end{equation*}

    Where $r_t(\theta) = \frac{\pi_\theta(a_t|s_t)}{\pi_{\theta}^{\text{old}}(a_t|s_t)}$ (probability ratio) and $\hat{A}_t$ is the advantage.
    \pause


\end{frame}

% ----------------------------------------------------------------
\section{RLHF: The Classic Approach}
% ----------------------------------------------------------------

\begin{frame}{Reinforcement Learning from Human Feedback (RLHF)}
    RLHF (popularized by OpenAI's InstructGPT, Ouyang et al., 2022) adapts RL principles to LLM training.
    \pause

    \begin{itemize}
        \item \textbf{Agent (Policy $\pi$):} The Language Model.
        \item \textbf{Action:} Generating the next token.
        \item \textbf{State:} The sequence of tokens generated so far (context).
        \item \textbf{Reward:} How do we score the quality of a complete response?
    \end{itemize}
    \pause
    \begin{alertblock}{The Challenge}
        We cannot mathematically define a perfect reward function for human language (e.g., a ``helpfulness score'').
    \end{alertblock}
    \pause
    \begin{block}{The Solution}
        We train a separate \textbf{Reward Model} (RM) to act as a proxy for human judgment.
    \end{block}
\end{frame}

\begin{frame}{The RLHF Pipeline: Three Stages}
    RLHF involves three distinct stages.

    % [Third-party paper figure omitted from the public snapshot; consult the cited primary paper.]
    \pause

    \begin{enumerate}
        \item Supervised Fine-Tuning (SFT) (Warm Start)
        \item Reward Model (RM) Training
        \item RL Optimization (PPO)
    \end{enumerate}
\end{frame}

\begin{frame}{}
\textbf{Stage 1: SFT (The Baseline)}
    \begin{block}{Goal: Establish a capable starting point.}
       \pause
    \begin{itemize}
        \item We fine-tune the base model on high-quality instruction-response pairs.
        \pause
        \item This provides a strong initial policy ($\pi_{\text{SFT}}$) that understands instructions.
        \pause
        \item Necessary because starting RL directly on a base model is highly unstable; the exploration space is too vast.
    \end{itemize}
    \end{block}

\textbf{Stage 2: Training the Reward Model (RM)}
    \begin{block}{Goal: Create a model $r_\phi(x, y)$ that outputs a scalar score}
    \pause

    \textbf{1. Preference Data Collection:}
    \begin{itemize}
        \item A prompt $x$ is sampled.
        \pause
        \item The SFT model generates multiple outputs $(y_1, y_2, ..., y_k)$.
        \pause
        \item Human labelers rank these outputs from best to worst.
        \pause
        \item This results in preference pairs: $(x, y_w, y_l)$, where $y_w$ (winner) is preferred over $y_l$ (loser).
    \end{itemize}
    \pause

    \textbf{2. Model Architecture:} Typically initialized from the SFT model with the final layer replaced by a regression head.
\end{block}


\end{frame}

\begin{frame}{Stage 2: The RM Loss Function}
    We use the \textbf{Bradley-Terry model}, which assumes the probability of preferring $y_w$ over $y_l$ depends on the difference in their underlying latent scores (rewards):
    \pause

    $$ P(y_w \succ y_l | x) = \sigma(r_\phi(x, y_w) - r_\phi(x, y_l)) $$
    \pause

    Where $\sigma$ is the sigmoid function. The RM is trained by minimizing the negative log-likelihood (a standard classification loss):
    \pause

    $$ \mathcal{L}_{\text{RM}}(\phi) = -\mathbb{E}_{(x, y_w, y_l) \sim D} [\log \sigma(r_\phi(x, y_w) - r_\phi(x, y_l))] $$
    \pause

    \begin{block}{Intuition}
        We maximize the margin between the score of the preferred response and the score of the rejected response.
    \end{block}
\end{frame}

\begin{frame}{Stage 3: RL Optimization (PPO)}
    \begin{block}{Goal: Optimize the LLM policy $\pi_\theta$ to generate responses that maximize the reward from the trained RM $r_\phi$.}

    \pause

    \begin{itemize}
        \item The policy $\pi_\theta$ (initialized from $\pi_{\text{SFT}}$) generates responses.
        \pause
        \item The (now frozen) RM scores these responses.
        \pause
        \item The policy weights $\theta$ are updated using an RL algorithm.
        \pause
        \item \textbf{Algorithm:} This lecture develops the stage using \textbf{Proximal Policy Optimization (PPO)}.
    \end{itemize}
    \end{block}
\end{frame}

\begin{frame}{Stage 3: An Additional Constraint Beyond PPO}
    \begin{alertblock}{The Risk: Reward Hacking}
        If we optimize purely for reward, the model might learn to generate nonsensical outputs that exploit weaknesses in the RM (Goodhart's Law). The RM is only a proxy.
    \end{alertblock}
    \pause

    \textbf{Two Levels of Constraints:}
    \pause
    \begin{enumerate}
        \item \textbf{PPO's built-in trust region:} Keeps $\pi_{\theta}^{\text{new}}$ close to $\pi_{\theta}^{\text{old}}$ (step-to-step stability).
        \pause
        \item \textbf{KL penalty to reference:} Keeps $\pi_{\theta}$ close to $\pi_{\text{ref}} = \pi_{\text{SFT}}$ (prevents drift from coherent language).
    \end{enumerate}
    \pause

    \begin{block}{The Complete RLHF Optimization Objective}
    $ \text{maximize } \mathbb{E}_{x, y \sim \pi_\theta} \left[ r_\phi(x, y) - \beta \cdot D_{KL}(\pi_\theta(y|x) || \pi_{\text{ref}}(y|x)) \right] $
    \end{block}
    \pause

    $\beta$ controls the trade-off: maximize reward while staying close to the original SFT model throughout training.
    \pause

    \textbf{Why both?} PPO prevents catastrophic step-wise updates; KL prevents long-term reward hacking.
\end{frame}

\begin{frame}{RLHF and PEFT: Practical Considerations}
    \begin{block}{
        Training large models using RLHF (PPO) is extremely memory intensive.}

    \pause
    During the PPO phase, we must load multiple models simultaneously:
    \begin{enumerate}
        \item The active policy $\pi_\theta$ (requires gradients).
        \pause
        \item The reference policy $\pi_{\text{ref}}$ (for KL calculation).
        \pause
        \item The Reward Model $r_\phi$.
        \pause
        \item The Value Model/Critic (part of the PPO algorithm).
    \end{enumerate}
    \end{block}
    \pause
    \begin{alertblock}{Parameter-Efficient Training}
        The lecture proposed freezing base models and training low-rank adapters for
        the policy, reward model, and value model to reduce the trainable state.
        Product-specific memory and feasibility claims are omitted here and checked
        independently before publication.
    \end{alertblock}
\end{frame}

% ----------------------------------------------------------------
\section{Challenges and Alternatives: DPO}
% ----------------------------------------------------------------

\begin{frame}{The Drawbacks of RLHF/PPO}
    While effective, the standard RLHF pipeline is notoriously difficult to implement:
    \pause

    \begin{itemize}
        \item \textbf{Complexity:} Requires training multiple models sequentially and managing a complex RL pipeline.
        \pause
        \item \textbf{Instability:} PPO is highly sensitive to hyperparameters and difficult to debug.
        \pause
        \item \textbf{Computational Cost:} High memory usage and slow sampling during the RL phase.
    \end{itemize}

    \vspace{0.5cm}
    \pause
    \begin{block}{So ..}
        Can we optimize the policy using preference data \textit{without} explicit Reward Modeling and complex Reinforcement Learning?
    \end{block}
\end{frame}

\begin{frame}{Direct Preference Optimization (DPO)}
    \textbf{Direct Preference Optimization (DPO)} (Rafailov et al., 2023) provides an elegant, stable, and efficient alternative.
    \pause

    \begin{block}{The ``Aha!'' Moment}
        In RLHF, the optimal policy $\pi^*$ that maximizes reward $r$ while staying close to $\pi_{\text{ref}}$ has a closed-form solution:
        $ \pi^*(y|x) \propto \pi_{\text{ref}}(y|x) \cdot \exp\left(\frac{r(x,y)}{\beta}\right) $
    \end{block}
    \pause

    We can \textbf{invert} this to express the reward in terms of the policy:
    $ r(x,y) = \beta \log \frac{\pi^*(y|x)}{\pi_{\text{ref}}(y|x)} + \text{constant} $
    \pause

    \begin{alertblock}{The Key Insight: Off-Policy Connection}
    This is the \textbf{log importance sampling ratio} from off-policy RL!
    \begin{itemize}
        \item $\pi_{\text{ref}}$ (SFT model) = \textbf{behavior policy} (generated the data)
        \item $\pi_\theta$ = \textbf{target policy} (what we're optimizing)
        \item The reward captures how much the target deviates from the reference
    \end{itemize}
    \end{alertblock}
    \pause

    \textbf{Implication:} We don't need to learn $r$ explicitly---just optimize the importance ratio directly!
\end{frame}

\begin{frame}{DPO: Cutting Out the Middleman}
    \textbf{The DPO Strategy:}
    \pause
    \begin{enumerate}
        \item Substitute $r(x,y) = \beta \log \frac{\pi_\theta(y|x)}{\pi_{\text{ref}}(y|x)}$ into the Bradley-Terry preference model.
        \pause
        \item This gives us a loss function that depends \textit{only} on $\pi_\theta$ and the preference data.
        \pause
        \item Optimize $\pi_\theta$ directly to maximize preference classification accuracy.
    \end{enumerate}
    \pause

    \begin{block}{Analogy}
        \textbf{RLHF:} Train a judge (RM) $\rightarrow$ Then try to please the judge (PPO) \\[0.2cm]
        \textbf{DPO:} Directly learn what pleases people by comparing examples
    \end{block}
    \pause

    \textbf{Result:} DPO treats alignment as a \textbf{classification problem}---given $(x, y_w, y_l)$, increase probability of $y_w$ and decrease probability of $y_l$, weighted by the policy ratio.
\end{frame}

\begin{frame}{The DPO Loss Function}
    DPO maximizes the relative log probability of preferred completions over dispreferred completions, inherently maintaining the KL constraint.
    \pause

    \begin{equation*}
    \begin{split}
        \mathcal{L}_{DPO}(\pi_\theta; \pi_{\text{ref}}) = -\mathbb{E}_{(x, y_w, y_l) \sim D} \Bigg[ \log \sigma \Bigg( &\beta \log \frac{\pi_\theta(y_w|x)}{\pi_{\text{ref}}(y_w|x)} \\
        &- \beta \log \frac{\pi_\theta(y_l|x)}{\pi_{\text{ref}}(y_l|x)} \Bigg) \Bigg]
    \end{split}
    \end{equation*}
    \pause

    \begin{block}{Intuition}
    We increase the likelihood of $y_w$ and decrease the likelihood of $y_l$. The degree of change is weighted by how much the current policy already diverges from the reference policy (controlled by $\beta$).
    \end{block}
\end{frame}

\begin{frame}{RLHF vs DPO: Sequential vs End-to-End}
    \textbf{The Fundamental Difference:}
    \pause

    \begin{columns}[T]
    \column{0.48\textwidth}
    \begin{block}{RLHF: Two-Stage}
        \textbf{Stage 1:} Train RM on preferences
        $ r_\phi(x,y) \leftarrow \text{preference data} $
        \textbf{Then freeze} $r_\phi$ \\[0.3cm]
        \textbf{Stage 2:} Optimize policy
        $ \pi_\theta \leftarrow \text{maximize } r_\phi(x,y) $
    \end{block}
    \pause

    \textcolor{red}{\textbf{Issue:}} No gradient flow between stages!
    \begin{itemize}
        \item RM might be wrong
        \item Distribution shift
        \item Can't correct RM errors
    \end{itemize}

    \column{0.48\textwidth}
    \pause
    \begin{block}{DPO: End-to-End}
        \textbf{Single stage:} Directly optimize policy on preference data
        $ \pi_\theta \leftarrow \text{preference data} $
        \\[0.3cm]
        Reward function is \textit{implicit}---derived analytically from $\pi_\theta$ and $\pi_{\text{ref}}$
    \end{block}
    \pause

    \textcolor{green!50!black}{\textbf{Benefit:}} Full gradient flow!
    \begin{itemize}
        \item Policy learns from preferences directly
        \item No frozen approximation
        \item Simpler optimization
    \end{itemize}
    \end{columns}
    \pause

    \vspace{0.3cm}
    \begin{alertblock}{Key Insight}
        DPO eliminates the ``frozen reward model bottleneck'' by optimizing end-to-end.
    \end{alertblock}
\end{frame}

\begin{frame}{DPO Advantages: Pipeline Comparison}
    \begin{center}
% [Third-party paper figure omitted from the public snapshot; consult the cited primary paper.]
\end{center}
    \pause

    \textbf{Key Advantages:}
    \begin{itemize}
        \item \textbf{Simplicity:} No PPO, no explicit RM, no sampling during training.
        \item \textbf{Stability:} Standard supervised learning optimization.
        \item \textbf{Efficiency:} Faster and lower VRAM usage.
        % [Comparative performance claim omitted; the book checks empirical claims against primary results.]
    \end{itemize}
\end{frame}

% ----------------------------------------------------------------
\section{Modern Efficient Methods}
% ----------------------------------------------------------------

\begin{frame}{Group Relative Policy Optimization (GRPO)}
    \textbf{GRPO}, as presented in the 2024 lecture notes, further simplifies RL-based alignment by removing the critic/value model.
    \pause

    \textbf{Key Innovation:} Group-based advantage estimation
    \begin{itemize}
        \item Sample multiple outputs $(y_1, ..., y_G)$ for each prompt $x$.
        \pause
        \item Compute rewards $r_i$ for each output using the RM.
        \pause
        \item Estimate advantage relative to the group mean:
        $$ \hat{A}_i = r_i - \frac{1}{G}\sum_{j=1}^G r_j $$
        \pause
        \item No separate value network needed!
    \end{itemize}
    \pause

    \textbf{Why is GRPO Efficient?}
    \begin{itemize}
        \item Eliminates one large model from memory (the critic).
        \pause
        % [Comparative stability claim omitted; the book checks empirical claims against primary results.]
        \pause
        % [Product-current deployment claim omitted from the public snapshot.]
    \end{itemize}
\end{frame}

\begin{frame}{Why Choose GRPO over DPO?}
    \textbf{The Key Difference: Online vs Offline Learning}
    \pause

    \begin{columns}[T]
    \column{0.48\textwidth}
    \begin{block}{DPO (Offline)}
        \begin{itemize}
            \item Learns from \textit{fixed} preference dataset
            \pause
            \item Preferences collected from SFT model outputs
            \pause
            \item Cannot explore beyond training data
            \pause
            \item \textcolor{red}{Distribution mismatch:} training data is static
        \end{itemize}
    \end{block}

    \column{0.48\textwidth}
    \pause
    \begin{block}{GRPO (Online)}
        \begin{itemize}
            \item Generates \textit{new} samples during training
            \pause
            \item Gets immediate reward feedback
            \pause
            \item Can discover better solutions through exploration
            \pause
            \item \textcolor{green!50!black}{Always on-policy:} learns from its own generations
        \end{itemize}
    \end{block}
    \end{columns}
    \pause

    \vspace{0.3cm}
    \textbf{When to use GRPO?}
    \begin{itemize}
        \item Complex reasoning tasks where exploration matters
        \item When you want the model to go beyond the training data
        \item When you have a good reward model and can afford online sampling
    \end{itemize}
\end{frame}



\begin{frame}{Two Diverging Paths in Alignment}

    \begin{columns}[T]
    \column{0.48\textwidth}
    \begin{block}{Path 1: Offline Simplification}
        \textbf{For:} Instruction following, safety, general helpfulness
        \pause
        \vspace{0.3cm}
        \textbf{Method:} Supervised-like approaches (DPO, KTO (unary feedback), IPO, Rejection Sampling + SFT)
        \pause

        \vspace{0.3cm}

        \textbf{Why:}
        \begin{itemize}
            \item Fixed preference data sufficient
            \item Simpler, more stable
            \item No need for exploration
        \end{itemize}
    \end{block}

    \column{0.48\textwidth}
    \pause
    \begin{block}{Path 2: Better Online RL}
        \textbf{For:} Complex reasoning, math, code generation
        \pause

        \vspace{0.3cm}
        \textbf{Method:} Improved RL algorithms (PPO, GRPO, Outcome-based RM)
        \pause
        \vspace{0.3cm}

        \textbf{Why:}
        \begin{itemize}
            \item Need exploration/search
            \item Model must discover solutions
            \item Verifiable rewards (unit tests, proofs)
        \end{itemize}
    \end{block}
    \end{columns}
    \pause

    \vspace{0.3cm}
    \begin{alertblock}{The Principle}
        Choose your alignment method based on whether your task benefits from \textbf{exploration} (RL) or just needs \textbf{preference fitting} (offline/supervised).
    \end{alertblock}
\end{frame}

% ----------------------------------------------------------------
\section{Future Directions and Summary}
% ----------------------------------------------------------------

\begin{frame}{Evolving Methods in Alignment}
    The lecture surveyed additional approaches beyond RLHF and DPO:
    \pause

    \begin{itemize}
        \item \textbf{KTO (Kahneman-Tversky Optimization):} (Ethayarajh et al., 2024). Argues that pairwise preferences (A>B) are unnecessary. KTO aligns models using only unary feedback (is this output ``good'' or ``bad''?). This drastically simplifies data collection.
        \pause
        \item \textbf{IPO (Identity Preference Optimization):} Addresses theoretical issues in DPO regarding long training runs by adding regularization.
        \pause
        \item \textbf{RLAIF (RL from AI Feedback):} Using another language model to generate preference labels instead of humans.
        \begin{itemize}
            % [Product-current cost and scalability claim omitted from the public snapshot.]
            \item Cons: Distills the preferences and biases of the teacher model.
        \end{itemize}
        \pause
        \item \textbf{Constitutional AI (CAI):} (Anthropic). RLAIF guided by a specific set of written principles (a ``constitution'').
    \end{itemize}
\end{frame}

\begin{frame}{Practical Tools}
    % [Product-current framework and platform recommendations omitted from this public snapshot.]
    The original lecture included an ecosystem survey here. Current product and API
    recommendations are checked separately before publication in the book.
\end{frame}

\begin{frame}{Summary: The Alignment Journey}
    \begin{enumerate}
        \item \textbf{Pre-training:} Develops core capabilities.
        \pause
        \item \textbf{SFT:} Teaches instruction following and format.
        \pause
        \item \textbf{Preference Tuning:} Aligns behavior with human values (HHH).
    \end{enumerate}

    \pause
    \begin{block}{Key Takeaways}
        \begin{itemize}
            \item Alignment is crucial for making LLMs useful and safe.
            \pause
            \item RLHF (PPO) was the foundational approach but is complex and costly.
            \pause
            \item DPO offers a simpler, stable, and efficient alternative by directly optimizing preferences.
            \pause
            \item The lecture connects PEFT (LoRA) to the resource demands of alignment training.
        \end{itemize}
    \end{block}
    \tiny
\textbf{References and Further Reading}
    \begin{itemize}
        \item Ouyang, L., et al. (2022). Training language models to follow instructions with human feedback. (InstructGPT) \href{https://arxiv.org/abs/2203.02155}{arXiv:2203.02155}

        \item Schulman, J., et al. (2017). Proximal Policy Optimization Algorithms. \href{https://arxiv.org/abs/1707.06347}{arXiv:1707.06347}

        \item Rafailov, R., et al. (2023). Direct Preference Optimization: Your Language Model is Secretly a Reward Model. \href{https://arxiv.org/abs/2305.18290}{arXiv:2305.18290}

        \item Ethayarajh, K., et al. (2024). KTO: Model Alignment as Prospect Theoretic Optimization. \href{https://arxiv.org/abs/2402.01306}{arXiv:2402.01306}

        \item Bai, Y., et al. (2022). Constitutional AI: Harmlessness from AI Feedback. \href{https://arxiv.org/abs/2212.08073}{arXiv:2212.08073}

        \item Askell, A., et al. (2021). A General Language Assistant as a Laboratory for Alignment. (The Anthropic HHH paper). \href{https://arxiv.org/abs/2112.00861}{arXiv:2112.00861}
    \end{itemize}
\end{frame}

\end{document}
